\chapter{Methodology}
\label{ch:methodology}

This chapter describes the experimental platform in the detail needed to reproduce the study. It sets out the diagnostic task and why it was chosen, the five energy functions, the two samplers and the factorial space of configurations in which they were run, the evaluation metrics and the reasoning behind them, the GFlowNet training setup, the constrained-generation extension, and finally the four diagnostic experiments that were designed specifically to isolate the mechanisms behind the main results.

\section{Task: Masked Token Recovery}
\label{sec:meth-task}

A sampler can only be studied if there is a way to tell whether it is being guided toward good sequences or merely wandering. This requires a task with a known target. We therefore adopt \textbf{masked token recovery}, also called infilling, as the diagnostic bench. A sequence is drawn from the corpus, one or more of its tokens are replaced with random tokens, and the sampler is asked to restore a coherent sequence by resampling only the corrupted positions while the remainder is held fixed. Because the original sequence is known, the recovery can be evaluated, and because the corruption and the context are controlled, the behaviour of the sampler can be attributed to the method rather than to the difficulty of an open-ended generation.

An important subtlety shapes the choice of evaluation metric and is worth stating at the outset. The token a sampler recovers need not be the original token in order to be correct. If the corrupted word \emph{roommate} is recovered as \emph{dorm}, the sequence is perfectly coherent even though it does not match the reference. An exact-match criterion would wrongly score this as a failure. For this reason the primary evaluation is based not on token identity but on how well the recovered sequence fits its context under the model's own predictive distribution, as described in Section~\ref{sec:meth-metrics}. The corpus is \textbf{ROCStories} \citep{mostafazadeh2016corpus}, a collection of short five-sentence everyday narratives. It is chosen for three reasons that matter for a controlled diagnostic: its sentences are short and syntactically simple, so a single corrupted token has a well-defined and locally checkable correct recovery; its domain is narrow and consistent, so the model's uncertainty reflects the sampler's behaviour rather than the open-endedness of the text; and it is the corpus used by the amortized-inference reference implementation of \citet{hu2024amortizing}, so the GFlowNet half of the study operates on the same distribution as the work it builds on. Masked recovery, rather than open-ended generation, is chosen as the diagnostic task precisely because it supplies a known target: a sampler can only be judged to be guiding well or poorly if there is a ground truth against which its output can be scored, and open-ended generation provides none.

\section{Energy Functions}
\label{sec:meth-models}

The study evaluates five energy functions. The reference energy is a \textbf{GPT-2 Large} model \citep{radford2019language}, of $774$ million parameters and embedding dimension $D = 1280$, supervised fine-tuned on ROCStories; this same fine-tuned model is used throughout as the base for the amortized experiments, so that the Langevin and GFlowNet halves of the study share weights and are not confounded by a difference in the underlying model. Three \textbf{GFlowNet-tuned variants} of this base model are included, obtained as described in Section~\ref{sec:meth-gfn}: two trained with the unnormalized length setting, at $500$ and $2000$ optimization steps respectively, and one trained with the length-normalized setting at $500$ steps. Finally, a \textbf{Llama-3 8B} model \citep{dubey2024llama}, with embedding dimension $D = 4096$, serves as a cross-architecture control. Because Llama-3 uses a different tokenizer and a different embedding geometry, and consequently a different step-size scale, it is never placed on a shared numerical axis with the GPT-2 family; it is used to establish that qualitative findings hold across architectures. The four GPT-2-based models share the tokenizer, the corruption seeds, and the step-size schedule, and their numbers are therefore directly comparable without qualification.

\section{Samplers and Configurations}
\label{sec:meth-configs}

The two samplers of Section~\ref{sec:bg-samplers}, the Discrete Langevin Sampler (DLS) and the Continuous Langevin Sampler (CLS), are each run under a factorial combination of design choices. The two are studied together, rather than either alone, because they represent the two opposite strategies for confronting the discreteness of text and therefore isolate different failure modes: the DLS, following \citet{grathwohl2021oops} and \citet{zhang2022langevin}, keeps the state in token space and uses the gradient only to build a proposal, which tests whether the gradient is an informative ranking signal; the CLS, following \citet{kumar2022gradient}, relaxes into continuous space and runs genuine Langevin dynamics, which tests whether continuous Markov chain Monte Carlo is compatible with a projected discrete energy. A finding that holds for both is a finding about the energy rather than about one sampler's idiosyncrasies, which is why both are needed. The \textbf{proposal method} is one of three: \emph{policy}, which uses the true gradient of the energy; \emph{grad-norm-preserved random}, which replaces the gradient with a random vector rescaled to the exact norm of the true gradient; and \emph{random}, which uses a random vector of random magnitude. The comparison between these three is the central ablation of the thesis, isolating respectively the full gradient, the effect of its direction, and a pure baseline. The \textbf{Metropolis--Hastings correction} is either enabled or disabled. \textbf{Gradient normalization}, which rescales the gradient to unit norm for numerical stability, is either enabled or disabled; this choice interacts with the proposal-method ablation in a way that turns out to be critical, as explained in Section~\ref{sec:results-fallacy}. The \textbf{annealing schedule} runs for either $50$ or $100$ steps, with the step size decreasing linearly; for the GPT-2 family the schedule runs from $\epsilon_{\text{start}} = 10.5$ to $\epsilon_{\text{end}} = 0.1$, a scale established by the calibration described in Section~\ref{sec:results-stepsize}. An \emph{oracle} variant, in which the ideal step size at each iteration is selected using knowledge of the ground-truth token, is included as an upper-bound reference; its purpose is to rule out the possibility that the samplers fail merely for want of a well-tuned schedule, since if even an oracle-tuned schedule does not let the gradient outperform a random direction, the failure cannot be one of calibration. In total, across the five energy functions, the study comprises $145$ sampling configurations, each evaluated on $200$ sequences.

\section{Evaluation Metrics}
\label{sec:meth-metrics}

Three quantities are recorded. The first is the \textbf{Euclidean distance} in embedding space between the recovered token and the ground-truth token. This is the most obvious measure of recovery, but it proves to be unreliable, for a reason established quantitatively in Section~\ref{sec:results-aniso}: because language-model embeddings are anisotropic, occupying a narrow cone, Euclidean distance compresses the distinction between good and bad tokens, and a token can be geometrically close to the target while being grammatically inadmissible. It is reported for completeness and as an object of study in its own right, not relied upon as the primary criterion.

The primary criterion is the \textbf{KL divergence} between the model's predictive distribution at the recovered position and a reference distribution, which measures how well the recovered token fits its context, that is, its \emph{contextual coherence}. A recovered sequence that the model finds locally natural has low divergence regardless of whether it matches the reference token, which is the property required of a metric on this task. We note one limitation of the divergence measure, which is that its magnitude depends on the position of the corrupted token within the sequence: corrupting an early token perturbs the conditional probability of a long suffix and yields a large divergence, whereas corrupting a final token yields almost none. For this reason no model is ranked on small differences in KL divergence; the central results are argued from the \emph{absence} of a gap between conditions rather than from a small measured gap, which is robust to this limitation in a way that a small measured difference would not be. The KL divergence \citep{kullback1951information} is chosen as the primary metric over exact-match accuracy because, as noted above, a recovered token that differs from the reference may still be perfectly coherent, and a metric that penalized such recoveries would confound the quality of the sampler with the arbitrariness of the reference; a divergence between predictive distributions measures contextual fit directly and is insensitive to which particular admissible token was chosen. Finally, for the generation-quality analyses, we record \textbf{repetition rate}, the fraction of repeated $n$-grams, and \textbf{distinct-$n$}, which are the standard measures of the degeneration studied by \citet{holtzman2020curious} and are therefore the appropriate instruments for detecting the likelihood trap; and for the constrained-generation extension we record classifier accuracy and \textbf{BERTScore} \citep{zhang2020bertscore}, the latter chosen because it measures semantic similarity through contextual embeddings rather than surface overlap and so does not penalize valid paraphrases of the reference.

\section{GFlowNet Training}
\label{sec:meth-gfn}

The GFlowNet is chosen as the amortized alternative, over reward-maximizing reinforcement learning, for a reason specific to the question of this thesis: a reward-maximizing agent would converge on the single highest-reward sequence and so could not distinguish a sampler that explores the energy from one that merely finds its mode, whereas a GFlowNet \citep{bengio2021flow, bengio2023gflownet} is trained to sample in proportion to reward and therefore tests amortized \emph{sampling} rather than amortized optimization. Crucially for the argument, it only ever evaluates the reward and never differentiates it, so it is blind to exactly the gradient pathology that afflicts the Langevin samplers, which makes it the ideal instrument for asking whether that pathology can be sidestepped. The GFlowNet-tuned variants are produced following the amortized-inference formulation of \citet{hu2024amortizing}, using their modified sub-trajectory-balance objective \citep{malkin2022trajectory, madan2023learning}. The policy is a low-rank adapter \citep{hu2022lora} on the frozen GPT-2 Large base, and the reward is the log-likelihood assigned by that frozen base to the completed sequence, so that the GFlowNet is trained to sample in proportion to the model's own likelihood. The task is a story-infilling formulation in which, given a preceding context $X$ and a known ending $Y$, the policy generates a connecting span $Z$, and the reward scores the concatenation. A single hyperparameter, denoted \texttt{len\_beta}, controls length normalization of the reward: at $\texttt{len\_beta}=0$ the reward is the raw unnormalized log-likelihood, and at $\texttt{len\_beta}=1$ it is fully normalized by sequence length. The consequences of this single knob are one of the central findings of the GFlowNet study, developed in Section~\ref{sec:results-gfn}.

\section{Constrained Generation Extension}
\label{sec:meth-constrained}

To test the plug-and-play claim directly (RQ4), a constrained-generation extension augments the energy with a sentiment constraint. A sentiment classification head is trained on the representations of the GPT-2 Large base and used to construct the constraint term $C(\bm{x})$ in \eqref{eq:intro-energy}. Generation is then run under five constraint modes that isolate the contribution of the constraint gradient: \emph{lm\_only}, which uses only the fluency energy; \emph{cons\_only}, which uses only the constraint; \emph{full}, which combines them with weights $\beta_{\text{LM}} = 0.8$ and $\beta_{\text{C}} = 0.2$; \emph{random}, which replaces the combined gradient with noise; and \emph{cons\_random}, which replaces only the constraint gradient with noise while retaining the fluency gradient. The steering effect is measured as the change in the fraction of generations classified as the target sentiment relative to the unconstrained baseline. Both the discrete sampler in the formulation of this thesis and a MuCoLa-style continuous baseline are evaluated, on both infilling and continuation.

\section{Diagnostic Experiments}
\label{sec:meth-diagnostics}

The configurations above establish \emph{that} the methods fail. To establish \emph{why}, four additional diagnostic experiments were designed, each targeting one candidate mechanism. They do not sample; they measure properties of the energy landscape and of the gradient directly.

The \textbf{linearization experiment} tests whether the first-order Taylor surrogate that the discrete proposal uses to rank candidate tokens actually predicts the true change in sequence energy. For each of $200$ sequences and a masked position, the exact gradient at that position is computed, and for a stratified set of $2000$ candidate tokens the surrogate predicted change $\hat{\Delta}(v) = \vg^\top(\ve(v) - \ve(x_i))$ is compared against the true change $\Delta(v) = \log p(\bm{x}_{i\to v}) - \log p(\bm{x})$ obtained by an exact forward pass, and the embedding distance of each candidate is recorded so that the correlation can be examined as a function of distance. The true change is further decomposed into a \emph{self} term, the change in the log-probability of the candidate given its own left context, and a \emph{future} term, the change in the log-probability of the suffix, because the gradient with respect to the input embedding is structurally able to see only the latter.

The \textbf{acceptance experiment} tests the Metropolis--Hastings breakdown in the continuous sampler by logging, for every proposal, whether it crossed a Voronoi cell boundary, whether it was accepted, and the decomposition of the acceptance ratio into its target and proposal terms. The \textbf{trajectory experiment} logs the raw state of the sampler at every step for a small number of sequences, so that the path can be projected into two dimensions with both PCA and t-SNE, and it records the distance from the continuous state to the nearest token embedding, quantifying how far the sampler drifts from the token manifold. The \textbf{likelihood-trap experiment} decodes with several strategies, from greedy and beam search to ancestral sampling, and measures the relationship between the per-token likelihood of the generated text and its repetition rate, establishing on the study's own models whether the lowest-energy text is in fact the worst text. A final \textbf{anisotropy analysis} measures the distribution of inter-token distances and related geometric quantities in the two embedding spaces. All experiments use a resumable file-queue system with per-job VRAM gating, and all sampling code was verified against a reference implementation with a bitwise-equivalence test suite before the runs reported here were launched.
